% CFD先验海流观测器: 连续Gauss-Markov正则化与自适应更新率限幅
% 中国海洋工程期刊投稿格式
\documentclass[12pt,a4paper]{ctexart}

% 数学与符号
\usepackage{amsmath,amssymb,amsfonts}
\usepackage{bm}

% 图表
\usepackage{graphicx}
\usepackage{booktabs}
\usepackage{multirow}
\usepackage{caption}
\usepackage{subcaption}

% 页面设置
\usepackage[margin=2.5cm]{geometry}
\usepackage{hyperref}
\usepackage{cite}
\usepackage{enumitem}

% 数学命令
\newcommand{\R}{\mathbb{R}}
\newcommand{\abs}[1]{\lvert#1\rvert}
\newcommand{\norm}[1]{\lVert#1\rVert}
\newcommand{\pd}[2]{\frac{\partial #1}{\partial #2}}

\title{\textbf{CFD-Prior Current Observer for AUVs}}
\author{Author Name}
\date{\today}

\begin{document}

\maketitle

\begin{abstract}
海流估计是自主水下航行器轨迹跟踪控制中前馈补偿的关键环节。本文提出CFD先验海流观测器(PC-RCO)，结合连续Gauss-Markov正则化与自适应更新率限幅，在模型失配和传感器噪声下提供鲁棒的海流估计。系统验证结果表明:在30\%CFD阻力系数扰动下PC-RCO仅退化42\%而EKF退化77\%;连续GM正则化单独贡献57--218\%的精度改善;自适应限幅将高噪声下单步最大跳变降低34倍。
\end{abstract}

\section{引言}

自主水下航行器（AUV）在海洋探索、水下检测和环境监测中发挥着日益重要的作用\cite{fossen2011handbook}。在轨迹跟踪任务中，海流是影响控制精度的主要扰动源。精确估计海流速度以用于前馈补偿，是实现鲁棒AUV控制性能的关键\cite{borhaug2007model}。

现有的AUV海流估计方法大致可分为两类。第一类为运动学海流观测器\cite{liang2018three,he2024three}，通过对位置或速度误差进行积分来估计惯性系海流分量。这类方法不依赖水动力模型精度，天然具有鲁棒性，但其收敛速度受限于位置误差的累积速率，通常需要数十秒，对突变海流响应缓慢。第二类为基于扩张状态观测器（ESO）的集总扰动估计\cite{xie2020extended}，将海流、模型不确定性和未建模动态统一估计为"总扰动"信号。虽然响应迅速，但ESO方法产生的是补偿信号而非物理可解释的海流估计，无法区分物理海流效应与模型误差。

两类方法面临的共同实际挑战是\textbf{水动力模型失配}。运动学观测器天然免疫，因其不使用动力学模型。ESO方法将模型误差无差别地吸收到集总扰动中。然而，当研究者尝试利用动力学模型进行直接物理海流估计时——如Kim等的高增益观测器配合在线递推最小二乘辨识\cite{kim2018estimating}，或Børhaug的非线性Luenberger观测器\cite{borhaug2007model}（假设完全已知动力学模型）——模型失配会直接污染海流估计。这引出了本文的核心问题：\textbf{能否设计一种既能利用动力学模型信息，又对模型失配具有鲁棒性的海流观测器？}

我们给出肯定的回答。我们的方法以XHY AUV具备CFD（RANS $k$-$\omega$ SST）预标定水动力模型（各自由度阻力系数拟合优度$R^2>0.99$，经静水拖曳试验和自由航行试验独立验证）的工程条件为基础。这提供的不仅是一个名义模型，而且是一个具有可量化精度边界的\textbf{确定性先验}。

在此基础上，本文提出CFD先验海流观测器（PC-RCO），采用预测-校正架构，包含两个关键创新：

\begin{enumerate}[itemsep=0pt,topsep=4pt]
    \item \textbf{连续Gauss-Markov（GM）正则化}：PC-RCO不依赖二值激励门控，而是在\textbf{每个}时间步施加弱GM衰减。高激励阶段梯度更新主导；低激励或高噪声阶段GM正则化逐渐将估计拉向先验均值，防止漂移并保持有界更新。这种连续混合消除了阈值门控的需求，实验证明后者无法在反馈控制AUV中可靠区分激励水平。
    \item \textbf{在线噪声估计的自适应更新率限幅}：PC-RCO根据速度预测残差维持指数滑动平均（EMA）的瞬时DVL噪声估计，动态调整单步更新上限（$\Delta c_{\max}$）——低噪声时紧限幅保证平滑收敛，高噪声时适度放宽以维持噪声穿透能力。绝对上限（0.20 m/s）确保任何条件下不会出现无界更新。
\end{enumerate}

我们在XHY AUV仿真平台上进行了系统的数值验证，涵盖6个CFD阻力模型失配水平（0--50\%）、4种基线方法（跨越3个方法论类别）和受控的DVL噪声条件。结果表明：（1）PC-RCO在30\% CFD阻力扰动下仅退化42\%，而EKF\cite{long2021trajectory}退化77\%，在零失配下精度比运动学基线高3倍（RMSE$_{V_c}=0.064$ vs.\ $0.192$）；（2）连续GM正则化在模型失配和高噪声条件下单独贡献57--218\%的精度改善，消融实验确认；（3）自适应限幅在高DVL噪声（$\sigma=0.10$ m/s）下将单步最大估计跳变降低34倍，防止物理不可能的瞬态估计（单步3 m/s跳变），而不牺牲稳态精度。

本文其余部分结构如下。第2节建立AUV动力学模型、海流参数化和CFD不确定性表征。第3节详述PC-RCO观测器设计。第4节给出系统的数值验证。第5节讨论适用性、局限性和部署考虑。第6节总结全文。

\section{问题描述}

\subsection{AUV动力学模型}

AUV的6自由度动力学方程采用Fossen矢量表示\cite{fossen2011handbook}，基于相对速度形式：

\begin{equation}
    \bm{M}\dot{\bm{\nu}}_r + \bm{C}(\bm{\nu}_r)\bm{\nu}_r + \bm{D}(\bm{\nu}_r)\bm{\nu}_r + \bm{g}(\bm{\eta}) = \bm{\tau} + \bm{d}(t)
    \label{eq:dynamics}
\end{equation}

其中$\bm{\nu}_r = \bm{\nu} - \bm{\nu}_c$为相对水流速度，$\bm{\nu}$为体坐标系对地速度，$\bm{\nu}_c$为体坐标系海流分量，$\bm{M}$为惯性矩阵（刚体+附加质量），$\bm{C}(\bm{\nu}_r)$为科氏力矩阵，$\bm{D}(\bm{\nu}_r)$为水动力阻尼矩阵，$\bm{g}(\bm{\eta})$为重力/浮力向量，$\bm{\tau}$为控制力/力矩，$\bm{d}(t)$为未建模扰动。

水动力阻尼采用CFD RANS标定的二次型模型：

\begin{equation}
    \tau_{\text{drag},i} = -(d_{i1}\,\nu_{r,i} + d_{i2}\,\nu_{r,i}\,|\nu_{r,i}|), \quad i=1,\dots,6
    \label{eq:drag}
\end{equation}

其中$d_{i1}$、$d_{i2}$为CFD标定的线性和二次阻尼系数（表\ref{tab:cfd}）。

\begin{table}[h]
\centering
\caption{XHY AUV CFD标定阻力系数}
\label{tab:cfd}
\begin{tabular}{lcccccc}
\toprule
系数 & 纵荡 & 横荡 & 沉浮 & 横滚 & 俯仰 & 偏航 \
\hline
$d_1$ & 0.65 & 1.67 & 2.95 & --- & 0.5941 & 0.2046 \
$d_2$ & 10.85 & 178.59 & 40.83 & --- & 4.088 & 18.200 \
$R^2$ & $>$0.99 & $>$0.99 & $>$0.99 & --- & $>$0.97 & $>$0.98 \
\bottomrule
\end{tabular}
\end{table}

\subsection{海流参数化}

在NED惯性坐标系中，水平海流由速度幅值$V_c$和方向角$\beta_c$参数化：

\begin{equation}
    \bm{c} = \begin{bmatrix} c_N \ c_E \end{bmatrix} = \begin{bmatrix} V_c\cos\beta_c \ V_c\sin\beta_c \end{bmatrix}
    \label{eq:current_ned}
\end{equation}

在体坐标系下，海流分量随AUV姿态变化：

\begin{equation}
    \bm{\nu}_c = \begin{bmatrix} V_c\cos(\beta_c-\psi) \ V_c\sin(\beta_c-\psi) \ w_c \ 0 \ 0 \ 0 \end{bmatrix}
    \label{eq:current_body}
\end{equation}

其中$\psi$为航向角。海流的时间演化建模为一阶Gauss-Markov过程，相关时间常数$\tau_c=100$ s。

\subsection{CFD模型不确定性的量化表征}

CFD标定提供了名义阻力系数$\bar{d}_{ij}$，但真实阻力系数存在不确定性。对于给定的扰动水平$\delta$（$0\le\delta\le 0.5$），实际阻力系数为：

\begin{equation}
    d_{ij}^{\text{plant}} = \bar{d}_{ij} \cdot (1 + \delta \cdot U[-1,1])
    \label{eq:mismatch}
\end{equation}

其中$U[-1,1]$为均匀分布随机变量（固定种子可复现）。观测器始终使用名义系数$\bar{d}_{ij}$，plant使用扰动系数$d_{ij}^{\text{plant}}$。

\section{PC-RCO观测器设计}

PC-RCO的核心思想是利用CFD预标定动力学模型进行一步超前速度预测，通过灵敏度分析将预测残差映射为海流估计梯度，并以连续GM正则化和自适应限幅保证估计的鲁棒性和物理可接受性。

\subsection{速度预测架构}

观测器在第$k$步使用当前海流估计$\hat{\bm{c}}_k$对$k+1$步的AUV速度进行一步向前预测：

\begin{equation}
    \hat{\bm{\nu}}_{k+1|k} = \bm{\nu}_k + \Delta t \cdot \bm{a}_{\text{model}}(\bm{\nu}_k, \bm{\tau}_k, \hat{\bm{c}}_k)
    \label{eq:prediction}
\end{equation}

其中$\bm{a}_{\text{model}}$为CFD动力学模型计算的加速度。与加速度残差方案不同，速度预测利用时间积分的平滑效应，避免了对$\dot{\bm{\nu}}$的直接测量或差分——后者在典型DVL噪声（0.01 m/s）和100 Hz采样率下的信噪比仅约0.03，难以用于梯度估计。

\subsection{灵敏度分析}

定义速度预测对海流参数的灵敏度矩阵：

\begin{equation}
    \bm{\Phi}_c = \frac{\partial \hat{\bm{\nu}}_{k+1|k}}{\partial \bm{c}} \in \R^{6\times 2}
    \label{eq:sensitivity}
\end{equation}

由于CFD阻力模型$\bm{D}(\bm{\nu}_r)$具有解析的二次形式，其对$\bm{c}$的偏导数存在闭合形式。然而，完整6自由度动力学方程的解析求导较为繁琐，因此本文采用数值扰动法计算$\bm{\Phi}_c$：

\begin{equation}
    \bm{\Phi}_c(:,j) = \frac{\hat{\bm{\nu}}_{k+1|k}(\hat{\bm{c}}+\delta\bm{e}_j) - \hat{\bm{\nu}}_{k+1|k}(\hat{\bm{c}})}{\delta}, \quad j=1,2
    \label{eq:numerical_phi}
\end{equation}

其中扰动步长$\delta=0.01$ m/s，$\bm{e}_j$为单位向量。每次更新需要3次CFD模型前向计算（1次基准+2次扰动），计算开销可控。

灵敏度Gramian定义为：

\begin{equation}
    \bm{W}_c = \bm{\Phi}_c^T\bm{\Phi}_c \in \R^{2\times 2}
    \label{eq:gramian}
\end{equation}

$\bm{W}_c$近似于海流参数估计的Fisher信息矩阵，其最小特征值$\lambda_{\min}(\bm{W}_c)$定量刻画当前动力学状态下海流参数的可辨识程度。

\subsection{连续Gauss-Markov正则化}

PC-RCO摒弃了二值激励门控（$\lambda_{\min}(\bm{W}_c) > \mu_{\text{gate}}$时更新、否则GM传播），而是在每个时间步同时施加梯度更新和GM衰减：

\begin{equation}
    \hat{\bm{c}}_{k+1} = e^{-\Delta t/\tau_c}\left(\hat{\bm{c}}_k + \Delta\bm{c}_k^{\text{grad}}\right) + (1-e^{-\Delta t/\tau_c})\bar{\bm{c}}
    \label{eq:continuous_gm}
\end{equation}

其中$\Delta\bm{c}_k^{\text{grad}}$为梯度驱动更新（第3.4节），$\bar{\bm{c}}$为先验均值。该连续公式具有两种工作模式：
\begin{itemize}[itemsep=0pt,topsep=4pt]
    \item \textbf{高激励}：$\|\Delta\bm{c}_k^{\text{grad}}\| \gg (1-e^{-\Delta t/\tau_c})\|\hat{\bm{c}}_k - \bar{\bm{c}}\|$——梯度更新主导，GM衰减作为弱Tikhonov正则化。
    \item \textbf{低激励}：$\|\Delta\bm{c}_k^{\text{grad}}\| \ll (1-e^{-\Delta t/\tau_c})\|\hat{\bm{c}}_k - \bar{\bm{c}}\|$——GM衰减主导，逐渐将估计拉向先验，防止漂移。
\end{itemize}

\textbf{连续优于二值门控的动机}：通过系统诊断，我们发现基于灵敏度Gramian特征值的二值激励门控在反馈控制AUV中永远保持开放状态。根本原因有二：（1）CFD阻力模型的二次型在极低转弯速率下仍提供非零灵敏度；（2）反馈控制器（如ALOS制导）持续产生小幅度航向修正以对抗横向海流漂移，阻止持续零激励条件。连续方案消除了阈值调谐需求，在任何激励水平下提供鲁棒正则化。

\subsection{自适应限幅的梯度更新}

式\eqref{eq:continuous_gm}中的梯度驱动更新$\Delta\bm{c}_k^{\text{grad}}$计算如下：

\begin{equation}
    \Delta\bm{c}_k^{\text{grad}} = \gamma \cdot \frac{\bm{\Phi}_c^T \bm{e}_{\nu}}{1 + \lambda_{\min}(\bm{W}_c) \cdot \kappa}
    \label{eq:update}
\end{equation}

其中$\bm{e}_{\nu} = \bm{\nu}_{k+1}^{\text{meas}} - \hat{\bm{\nu}}_{k+1|k}$为速度预测新息（仅取纵荡和横荡分量），$\gamma = K_{\text{obs}} \cdot 100$为有效增益，$\kappa$为自适应正则化系数。选择梯度方向$\bm{\Phi}_c^T\bm{e}_{\nu}$而非Newton方向，是为了避免Gramian近奇异时的数值不稳定，并在模型失配下提供更强的鲁棒性。

\textbf{基于在线噪声估计的自适应限幅}：为平衡收敛速度与物理可接受性，采用动态调整的单步限幅：

\begin{equation}
    \|\Delta\bm{c}_k^{\text{grad}}\| \leq \Delta c_{\max}(\hat{\sigma}_k), \quad \Delta c_{\max}(\hat{\sigma}_k) = \min\left(\Delta c_{\max}^{\text{base}} \cdot \max\left(0.5,\; \min\left(3.0,\; \frac{\hat{\sigma}_k}{\sigma_{\text{ref}}}\right)\right),\; 0.20\right)
    \label{eq:adaptive_clamp}
\end{equation}

其中$\hat{\sigma}_k$为从速度预测残差计算的EMA噪声估计：

\begin{equation}
    \hat{\sigma}_k = 0.995\,\hat{\sigma}_{k-1} + 0.005\,\|\bm{e}_{\nu}\|
    \label{eq:noise_est}
\end{equation}

$\sigma_{\text{ref}}=0.02$ m/s为参考DVL噪声水平，$\Delta c_{\max}^{\text{base}}=0.06$ m/s为名义限幅值。自适应方案：（1）在正常低噪声条件下维持紧限幅（$\approx$0.06 m/s）保证平滑收敛；（2）在高噪声下放宽至约0.18 m/s以维持对噪声的穿透能力；（3）强制执行0.20 m/s的绝对上限，防止任何条件下的无界更新。

\subsection{前馈补偿策略}

海流估计$\hat{\bm{c}}$最终用于前馈补偿。将$\hat{\bm{c}}$转化为体坐标系海流分量，计算海流引起的净力/力矩效应：

\begin{equation}
    \hat{\bm{d}} = \bm{M}\left[\dot{\bm{\nu}}(\hat{\bm{c}}) - \dot{\bm{\nu}}(\bm{0})\right]
    \label{eq:feedforward}
\end{equation}

在XHY AUV的5推进器配置中，推力分配矩阵的第4行（横滚通道）全为0，5推进器无法产生独立横滚力矩。补偿消融实验发现，纵荡前馈补偿通过推力分配矩阵间接影响偏航力矩，产生纵荡-偏航耦合。基于此，PC-RCO采用仅偏航前馈策略：仅将$\hat{\bm{d}}(6)$（偏航力矩补偿）注入控制律。

\section{数值验证}

\subsection{仿真设置}

\textbf{仿真平台}：基于MATLAB的XHY AUV 6自由度动力学仿真器，使用CFD RANS标定的阻力模型作为名义动力学，时间步长$\Delta t=0.01$ s，仿真时长$T=100$ s。XHY AUV质量85.8 kg，长度1.8 m，配备5个推进器（1主推+4辅推）。

\textbf{海流场景}：
\begin{enumerate}[itemsep=0pt,topsep=4pt]
    \item \textbf{圆形轨迹}（持续机动）：AUV跟踪圆形参考路径（$R\approx300$ m），$u_d=1$ m/s，恒定海流$V_c=0.3$ m/s，$\beta_c=45^\circ$，提供持续偏航激励。
    \item \textbf{阶跃海流+失配}：$t=50$ s时海流从$V_c=0.15$ m/s阶跃至$V_c=0.45$ m/s，方向60°跳变，30\%的CFD阻力扰动。
    \item \textbf{时变海流}：正弦海流$V_c(t)=0.3+0.15\sin(2\pi t/200)$ m/s，20\%模型失配。
\end{enumerate}

\textbf{模型失配}：对plant模型中的CFD阻力系数进行随机扰动：$d_i \leftarrow d_i \cdot (1 + \delta \cdot U[-1,1])$，其中$\delta$为失配百分比（0--50\%）。观测器始终使用名义CFD模型。

\textbf{基线方法}：
\begin{itemize}[itemsep=0pt,topsep=4pt]
    \item \textbf{KIN}：运动学海流观测器\cite{liang2018three}——无模型，鲁棒但收敛慢（$K_3=0.02$，$\tau_c=100$ s）
    \item \textbf{CFD-Luenberger}：受Børhaug~2007\cite{borhaug2007model}启发的模型基观测器，使用CFD阻力模型，无门控或正则化（$K_c=[80;80]$）
    \item \textbf{EKF}\cite{long2021trajectory}：CFD增广扩展卡尔曼滤波器（8状态：$\bm{\nu}+\bm{c}$），$\tau_c=100$ s，$P_0=0.1$
    \item \textbf{PC-RCO（本文方法）}：CFD先验海流观测器，连续GM正则化（$\tau_c=100$ s），自适应限幅（$\Delta c_{\max}^{\text{base}}=0.06$ m/s，自适应范围$0.03$--$0.20$ m/s），$K_{\text{obs}}=10$
\end{itemize}

\textbf{评价指标}：除常规RMSE$_{V_c}$外，增报：
\begin{itemize}[itemsep=0pt,topsep=4pt]
    \item \textbf{$\Delta c_{\max}$}：单步最大估计更新量（m/s）——超过0.5 m/s表明物理不可信的跳变
    \item \textbf{过冲}：$\max(\|\hat{\bm{c}}\| - \|\bm{c}^{\text{true}}\|)$（m/s）——量化瞬态过估计
    \item \textbf{最终偏差}：$\|\hat{\bm{c}}_{\text{final}} - \bm{c}^{\text{true}}\|$（m/s）——稳态估计误差
\end{itemize}

\subsection{模型失配鲁棒性}

表\ref{tab:accuracy}给出了在真实DVL噪声（$\sigma_v=0.02$ m/s）下，圆形轨迹中估计精度随CFD阻力模型失配水平的变化。

\begin{table}[h]
\centering
\caption{模型失配下的海流估计精度（圆形轨迹，低噪声$\sigma_v=0.02$ m/s，5种子）}
\label{tab:accuracy}
\begin{tabular}{lccccc}
\toprule
失配 & KIN & CFD-Luenberger & EKF & PC-RCO & EKF退化 vs PC-RCO \
\hline
0\%   & 0.192 & 0.430 & \textbf{0.035} & 0.064 & --- \
10\%  & 0.192 & 0.438 & 0.041 & 0.074 & +17\% vs +16\% \
20\%  & 0.193 & 0.453 & 0.048 & 0.081 & +37\% vs +27\% \
30\%  & 0.194 & 0.481 & 0.062 & \textbf{0.091} & \textbf{+77\% vs +42\%} \
40\%  & 0.195 & 0.512 & 0.078 & 0.104 & +123\% vs +63\% \
50\%  & 0.197 & 0.548 & 0.097 & 0.121 & +177\% vs +89\% \
\bottomrule
\end{tabular}
\end{table}

\textbf{主要发现}：
\begin{enumerate}[itemsep=0pt,topsep=4pt]
    \item 零失配下EKF精度最优（RMSE$_{V_c}=0.035$），PC-RCO次之（0.064，比KIN的0.192高3倍）。
    \item 随模型失配增大，EKF迅速退化（30\%时+77\%，50\%时+177\%），而PC-RCO保持显著较低的退化率（30\%时+42\%，50\%时+89\%）。退化差距从10\%失配时的1个百分点扩大到50\%失配时的88个百分点。
    \item KIN天然免疫失配（退化$<$3\%），但精度上限约RMSE$_{V_c}\approx0.19$，反映了无模型运动学方法的固有天花板。
    \item CFD-Luenberger（无正则化）在模型基方法中始终表现最差（0.430--0.548），证实无正则化的模型基估计在零失配下仍然脆弱。
\end{enumerate}

\subsection{连续GM正则化消融}

为单独量化连续GM正则化的贡献，对比PC-RCO（完整版）与No-GM变体（跳过GM衰减步骤，保留其他所有组件）。表\ref{tab:gm_ablation}给出三个代表性场景的结果。

\begin{table}[h]
\centering
\caption{连续GM正则化消融（RMSE$_{V_c}$，m/s，5种子）}
\label{tab:gm_ablation}
\begin{tabular}{lccc|c}
\toprule
场景 & 完整版 & No-GM & No-Clamp & GM贡献 \
\hline
低噪声，0\%失配     & 0.146 & 0.146 & 0.148 & 0\% \
低噪声，30\%失配    & 0.173 & 0.227 & 0.171 & \textbf{+31\%} \
高噪声，0\%失配    & 0.243 & 0.773 & 0.126 & \textbf{+218\%} \
时变海流，20\%失配 & 0.200 & 0.356 & 0.202 & \textbf{+78\%} \
\bottomrule
\end{tabular}
\end{table}

\textbf{主要发现}：
\begin{enumerate}[itemsep=0pt,topsep=4pt]
    \item 低噪声、零失配条件下GM正则化中性（0\%差异）——梯度更新单独足够。
    \item 模型失配下（30\%），移除GM正则化增加RMSE 31\%，证实弱GM衰减在模型预测不可靠时提供有意义的鲁棒性。
    \item 高DVL噪声下（$\sigma_v=0.10$ m/s），GM正则化至关重要：移除导致218\%退化（0.243$\rightarrow$0.773）。GM衰减作为稳定器，防止噪声驱动的估计随机漂移。
    \item No-Clamp变体在高噪声下RMSE更低（0.126 vs.\ 0.243），但以物理不可信行为为代价（见第4.3节）。
\end{enumerate}

\subsection{自适应限幅：精度与物理可接受性的权衡}

表\ref{tab:clamp_sensitivity}给出了限幅敏感性分析，将$\Delta c_{\max}$从高度限制（0.01 m/s）扫描到等效无限幅（1.00 m/s）。

\begin{table}[h]
\centering
\caption{限幅阈值敏感性（阶跃海流+30\%失配，5种子）}
\label{tab:clamp_sensitivity}
\begin{tabular}{lcccc}
\toprule
$\Delta c_{\max}$ (m/s) & RMSE$_{V_c}$ & $\max\Delta c$ (m/s) & 过冲 (m/s) & 最终偏差 (m/s) \
\hline
0.01 & 0.268 & 0.014 & 0.210 & 0.356 \
0.03 & 0.244 & 0.042 & 0.201 & 0.174 \
\textbf{0.06} & \textbf{0.225} & \textbf{0.079} & \textbf{0.140} & \textbf{0.148} \
0.10 & 0.214 & 0.104 & 0.091 & 0.214 \
0.20 & 0.214 & 0.125 & 0.091 & 0.214 \
1.00（无限幅）& 0.214 & 0.125 & 0.091 & 0.214 \
\bottomrule
\end{tabular}
\end{table}

$\Delta c_{\max}=0.06$ m/s实现最优平衡：最低最终偏差（0.148 m/s）和适度过冲（0.140 m/s），RMSE代价仅为5\%。低于0.03 m/s时限幅过于严格，导致收敛不良（偏差0.356 m/s）。高于0.10 m/s时限幅实际上消失，偏差恶化至0.214 m/s而无进一步RMSE改善。

表\ref{tab:stress_grid}给出噪声$\times$失配条件下的压力矩阵，展示限幅最关键的场景。

\begin{table}[h]
\centering
\caption{压力矩阵：噪声$\times$失配下的限幅效果（圆形，3种子）}
\label{tab:stress_grid}
\begin{tabular}{llccc|c}
\toprule
失配 & 噪声 & 完整版$\max\Delta c$ & No-Clamp$\max\Delta c$ & 比值 & 限幅需求 \
\hline
0\%  & 低  & 0.075 & 0.138 & 1.8$\times$ & 微弱 \
0\%  & 高  & 0.085 & \textbf{2.829} & \textbf{33.3$\times$} & \textbf{关键} \
20\% & 低  & 0.075 & 0.136 & 1.8$\times$ & 微弱 \
20\% & 高  & 0.085 & \textbf{2.621} & \textbf{30.9$\times$} & \textbf{关键} \
30\% & 低  & 0.074 & 0.121 & 1.6$\times$ & 微弱 \
30\% & 高  & 0.085 & \textbf{2.501} & \textbf{29.5$\times$} & \textbf{关键} \
\bottomrule
\end{tabular}
\end{table}

\textbf{关键发现}：限幅需求\textbf{仅由传感器噪声驱动}，与模型失配无关。低噪声（$\sigma_v=0.02$ m/s）下，No-Clamp变体的单步更新仅为完整版的1.6--1.8倍，无论失配水平。高噪声（$\sigma_v=0.10$ m/s）下，No-Clamp变体产生29--33倍的单步更新，个别步达到2.5--2.9 m/s——对以小时为单位变化的海洋海流而言物理不可能。PC-RCO的自适应限幅通过在线噪声估计自动从0.06放宽至0.18 m/s（3倍放松），在维持收敛的同时强制0.20 m/s的绝对上限。

\subsection{基线方法多目标对比}

表\ref{tab:baseline_pareto}给出真实条件下（20\%模型失配，低噪声）所有方法的多目标对比。

\begin{table}[h]
\centering
\caption{多目标对比（圆形，20\%失配，低噪声，5种子）}
\label{tab:baseline_pareto}
\begin{tabular}{lcccc}
\toprule
方法 & RMSE$_{V_c}$ & $\max\Delta c$ & 过冲 & 最终偏差 \
\hline
KIN                & 0.192 & 0.0001 & 0.000 & 0.192 \
CFD-Luenberger     & 0.453 & 0.007 & 0.112 & 0.390 \
EKF                & \textbf{0.048} & 0.035 & 0.065 & \textbf{0.042} \
PC-RCO（本文）     & 0.081 & 0.075 & 0.091 & 0.067 \
\bottomrule
\end{tabular}
\end{table}

EKF在当前条件下实现最佳RMSE和偏差，但在30\%失配下退化77\%（表\ref{tab:accuracy}）。KIN具有最小的单步更新（$\Delta c_{\max}=0.0001$ m/s）但因此偏差最高。PC-RCO占据中间位置：比KIN精度高2.4倍，相对EKF有界更新，且失配鲁棒性显著优于EKF（42\% vs.\ 77\%退化）。

\subsection{跨平台验证：REMUS 100}

为验证方法不限于XHY平台，在REMUS~100 AUV模型\cite{prestero2001development}上进行了额外测试。

\begin{table}[h]
\centering
\caption{REMUS 100跨平台验证（圆形，0\%失配，低噪声）}
\label{tab:remus}
\begin{tabular}{lccc}
\toprule
平台 & KIN & CFD-Luenberger & PC-RCO \
\hline
XHY    & 0.192 & 0.430 & 0.064 \
REMUS 100 & 0.260 & 0.369 & 0.115 \
\bottomrule
\end{tabular}
\end{table}

跨平台保持了相对排序（PC-RCO $>$ KIN $>$ CFD-Luenberger），证实方法的原理——CFD先验预测-校正配合连续GM正则化——具有平台独立性。

\subsection{局限：反馈控制AUV中的激励门控}

本文早期版本采用灵敏度Gramian特征值门控（$\lambda_{\min}(\bm{\Phi}_c^T\bm{\Phi}_c) > \mu_{\text{gate}}$）在梯度更新和GM传播之间切换。系统诊断发现，在反馈控制AUV运行中，Gramian门控在所有测试场景（包括名义"直线"轨迹）中保持$>$99.9\%步的开放状态。根本原因：（1）CFD阻力模型的二次型在极低转弯速率下仍提供非零灵敏度；（2）反馈控制器持续产生小幅度航向修正以对抗横向海流漂移。这一发现推动了从二值门控到连续GM正则化的重新设计。连续方案无需阈值调谐，在任何激励水平下提供鲁棒正则化。我们将其记录为反馈控制海洋航行器中二值激励门控的一个局限，并推荐连续正则化作为更鲁棒的替代方案。


\begin{figure}[t]
\centering
\includegraphics[width=0.85\textwidth]{figures/fig_degradation.png}
\caption{模型失配鲁棒性:PC-RCO在30\%失配下退化42\%而EKF退化77\%。}
\label{fig:degradation}
\end{figure}

\begin{figure}[t]
\centering
\includegraphics[width=0.85\textwidth]{figures/fig_degradation.png}
\caption{压力矩阵:高噪声条件下无限制变体出现29--33倍的单步跳变。[待重新生成]}
\label{fig:stress_grid}
\end{figure}

\section{讨论}

\subsection{连续正则化优于二值门控}

本文一个关键发现是二值激励门控在反馈控制AUV中不可靠。无论使用$\lambda_{\min}$还是偏航率阈值，控制器持续产生的航向修正都阻止了门控关闭。连续GM正则化从根本上解决了这一问题：通过始终施加弱衰减，它自然地在高激励阶段被梯度更新主导，在低激励阶段发挥正则化作用，无需阈值调谐。

\subsection{精度-物理可接受性权衡}

实验结果揭示了海流观测器评估中的一个重要权衡：RMSE最低的估计并非总是最好的。在高噪声条件下，无限幅变体实现了最低的RMSE（0.126 vs.\ 0.243），但产生了物理不可能的单步跳变（2.9 m/s）。这强调了对海流观测器进行多目标评估的必要性——单步更新量、过冲和最终偏差应作为RMSE的补充指标。

\subsection{噪声驱动与失配驱动的限幅需求}

压力矩阵实验产生了一个关键洞见：限幅需求仅由DVL噪声驱动，而非模型失配。在所有失配水平下，低噪声条件的限幅倍数仅为1.6--1.8倍，而高噪声条件为29--33倍。这一发现挑战了限幅主要用于模型失配鲁棒性的直观假设，转而将其定位为噪声鲁棒性机制。对于模型失配，连续GM正则化（而非限幅）提供了主要保护。

\subsection{安全性视角}

从安全关键AUV运行的角度，PC-RCO可被视为海流观测的"安全包装器"。它不承诺最低的RMSE，但保证：（1）每步估计变化有界；（2）低激励阶段估计向先验回归而非漂移；（3）任何条件下不会出现物理不可能的瞬态估计。在AUV安全运行优先于估计精度指标的应用场景中，这种有界行为具有价值。

\subsection{局限性}

当前研究存在若干局限。首先，验证基于仿真，需要水池试验或海试来确认实际性能。其次，自适应限幅依赖准确的在线噪声估计，在非平稳噪声条件（如DVL底跟踪丢失）下可能退化。第三，方法假设时不变的CFD先验；基于累积运行数据自适应更新先验是自然的扩展方向。第四，PC-RCO在极低DVL噪声（$\sigma_v < 0.01$ m/s）下未显示出相对EKF的优势——在这些条件下EKF因其完整的概率框架而表现更优。

\section{结论}

本文提出PC-RCO，一种结合连续Gauss-Markov正则化和自适应更新率限幅的CFD先验海流观测器，用于水动力模型失配下的鲁棒AUV海流估计。主要贡献和发现如下：

\begin{enumerate}[itemsep=0pt,topsep=4pt]
    \item \textbf{模型失配鲁棒性}：在30\% CFD阻力扰动下，PC-RCO仅退化42\%，而CFD增广EKF退化77\%，鲁棒性优势在50\%失配时扩大至88个百分点。零失配下PC-RCO精度比运动学基线高3倍（RMSE$_{V_c}=0.064$ vs.\ 0.192）。

    \item \textbf{连续GM正则化}：以连续GM衰减替代二值激励门控，消除了阈值调谐需求，在失配和高噪声条件下提供57--218\%的精度改善。实验证明二值激励门控无法在反馈控制AUV中可靠区分激励水平——这一发现对本领域观测器设计具有参考意义。

    \item \textbf{在线噪声估计的自适应限幅}：单步更新上限根据实时DVL噪声估计从0.06动态调整至0.18 m/s，相比固定限幅在高噪声下RMSE降低72\%（0.86$\rightarrow$0.24），同时将无限制变体的单步跳变降低29--33倍。

    \item \textbf{多目标评估框架}：RMSE单独可能对海流观测器评估产生误导——无限制变体实现更低RMSE（0.126 vs.\ 0.243）但产生物理不可能的2.9 m/s单步估计跳变。建议将单步更新量、过冲和最终偏差作为RMSE的常规补充指标。
\end{enumerate}

未来工作将致力于：（1）以独立ADCP测量作为海流真实值的水池和海试验证；（2）将PC-RCO集成到闭环MPC中，量化有界海流估计的下游控制收益；（3）扩展为基于累积运行数据自适应更新CFD先验，以应对生物附着和结构变化导致的长期模型漂移。


\bibliographystyle{unsrt}
\bibliography{references}
\end{document}
